\newpage

\begin{quote}
\textbf{Versioning:} MinEnglish follows
\href{https://semver.org}{Semantic Versioning}.
\texttt{MAJOR.MINOR.PATCH} --- breaking rule changes increment MAJOR.
\texttt{0.x.y} indicates active pre-stable development; \texttt{1.0.0}
marks the first frozen specification.

\textbf{AI Attribution:} This specification was developed with extensive
assistance from large language models. \textbf{Claude Opus} (Anthropic),
\textbf{Gemini 2.5 Pro} (Google DeepMind), and \textbf{Antigravity}
(Google DeepMind) contributed to the drafting, formalization, revision,
and example generation documented herein. The intellectual framework,
design decisions, and authorship remain those of the human author.
\end{quote}

\begin{center}\rule{0.5\linewidth}{0.5pt}\end{center}

\newpage

\subsection{Abstract}\label{abstract}

MinEnglish is a formal simplification and clarification of Modern
English. Its aim is to \textbf{remove} the inconsistencies,
redundancies, and irregularities that burden English --- not to replace
it with something foreign. The result is a leaner, fully regular version
of English: same vocabulary, same phonemes, same word order --- but with
predictable spelling, explicit quantities, and precise time operators.

Because MinEnglish enforces an absolute mathematical structure with zero
exceptions, it yields two immediate practical benefits: it is
\textbf{Easy to Learn} (there are no irregular cases to memorize) and
\textbf{Easy to Type} (due to a strict lowercase axiom, dense phonetic
spelling, and single-character mathematical operators).

This specification documents the phonological, morphological, and
syntactic rules targeting the stable release, demonstrates a 35.0\%
average character reduction over standard English across a 35-sentence
corpus, and acknowledges the primary adoption challenge: human habit.

\begin{center}\rule{0.5\linewidth}{0.5pt}\end{center}

\newpage

\subsection{1. Key Features}\label{key-features}

\begin{quote}
\textbf{Mission:} MinEnglish simplifies, clarifies, and formalizes
Modern English. It removes constraints --- it does not add them. Every
rule below exists to \textbf{eliminate} an inconsistency or ambiguity
already present in English.
\end{quote}

{\def\LTcaptype{} % do not increment counter
\begin{longtable}[]{@{}
  >{\raggedright\arraybackslash}p{(\linewidth - 2\tabcolsep) * \real{0.1386}}
  >{\raggedright\arraybackslash}p{(\linewidth - 2\tabcolsep) * \real{0.8614}}@{}}
\toprule\noalign{}
\begin{minipage}[b]{\linewidth}\raggedright
Feature
\end{minipage} & \begin{minipage}[b]{\linewidth}\raggedright
What it removes from English
\end{minipage} \\
\midrule\noalign{}
\endhead
\bottomrule\noalign{}
\endlastfoot
\textbf{Easy to Learn} & Removes the concept of ``exceptions'' entirely.
There are no irregular verbs, unpredictable spellings, or confusing
contextual plurals. \\
\textbf{Easy to Type} & Removes excessive keystrokes and finger travel.
Relies on a strict lowercase axiom and mathematical operators to
prioritize digital efficiency. \\
\textbf{Phonetic Spelling} & Removes the burden of memorizing irregular
spelling. Every letter has one sound, every sound one letter.
(\texttt{knight} → \texttt{nait}, \texttt{enough} → \texttt{inuf}) \\
\textbf{Explicit Quantities} & Removes ambiguous pluralization.
\texttt{cats} vs \texttt{cat} is a noise-channel problem. \texttt{3cat}
is unambiguous. Singular is the default --- no suffix needed. \\
\textbf{Regular Time} & Removes 12+ irregular tense forms. One temporal
operator (\texttt{-1d}, \texttt{1d}, \texttt{:5Y}) does the work of
went/gone/had gone/will go/would have gone. \\
\textbf{Symbolic Operators} & Replaces multi-word constructs with single
tokens already familiar from math and logic (\texttt{\&}, ` \\
\textbf{Zero Exceptions} & Removes the entire category of irregular
verbs, irregular plurals, and silent letters. What works for one word
works for all words. \\
\end{longtable}
}

\begin{center}\rule{0.5\linewidth}{0.5pt}\end{center}

\newpage

\subsection{2. Phonology and
Orthography}\label{phonology-and-orthography}

\subsubsection{2.1 Vowels}\label{vowels}

To ensure acoustic discernibility across diverse linguistic backgrounds,
the vowel space is constrained to five cardinal points:

{\def\LTcaptype{} % do not increment counter
\begin{longtable}[]{@{}ll@{}}
\toprule\noalign{}
MinEnglish & English Equivalent \\
\midrule\noalign{}
\endhead
\bottomrule\noalign{}
\endlastfoot
\textbf{a} (/æ/) & c\textbf{a}t (cat) \\
\textbf{e} (/ɛ/) & b\textbf{e}d (bed) \\
\textbf{i} (/ɪ/) & b\textbf{i}t (bit) \\
\textbf{o} (/ɒ/) & h\textbf{o}t (hot) \\
\textbf{u} (/ʌ/) & b\textbf{u}t (but) \\
\end{longtable}
}

\subsubsection{2.2 Long Vowels and
Diphthongs}\label{long-vowels-and-diphthongs}

Diphthongs and elongated vowels are encoded as contiguous grapheme
pairs, preserving the 1-to-1 spelling mapping without requiring macron
diacritics.

{\def\LTcaptype{} % do not increment counter
\begin{longtable}[]{@{}ll@{}}
\toprule\noalign{}
MinEnglish & English Equivalent \\
\midrule\noalign{}
\endhead
\bottomrule\noalign{}
\endlastfoot
\textbf{ei} (/eɪ/) & d\textbf{ay}, m\textbf{a}ke (dei, meic) \\
\textbf{ii} (/iː/) & s\textbf{ee}, \textbf{ea}t (sii, iit) \\
\textbf{ai} (/aɪ/) & l\textbf{i}ke, m\textbf{y} (laic, mai) \\
\textbf{ou} (/oʊ/) & g\textbf{o}, n\textbf{o} (gou, nou) \\
\textbf{uu} (/uː/) & f\textbf{oo}d, t\textbf{oo} (fuud, tuu) \\
\textbf{au} (/aʊ/) & h\textbf{ow}, h\textbf{ou}se (hau, haus) \\
\textbf{oi} (/ɔɪ/) & b\textbf{oy} (boi) \\
\textbf{ar} (/ɑːr/) & c\textbf{ar} (car) \\
\textbf{or} (/ɔːr/) & f\textbf{or} (for) \\
\textbf{er} (/ɜːr/) & h\textbf{er}, b\textbf{ir}d (her, berd) \\
\end{longtable}
}

\subsubsection{2.3 Consonants}\label{consonants}

Consonantal mappings prioritize acoustic clarity while resolving
historical orthographic collisions (such as the hard-soft \emph{c}
boundary).

{\def\LTcaptype{} % do not increment counter
\begin{longtable}[]{@{}
  >{\raggedright\arraybackslash}p{(\linewidth - 6\tabcolsep) * \real{0.1119}}
  >{\raggedright\arraybackslash}p{(\linewidth - 6\tabcolsep) * \real{0.0746}}
  >{\raggedright\arraybackslash}p{(\linewidth - 6\tabcolsep) * \real{0.6269}}
  >{\raggedright\arraybackslash}p{(\linewidth - 6\tabcolsep) * \real{0.1866}}@{}}
\toprule\noalign{}
\begin{minipage}[b]{\linewidth}\raggedright
Grapheme
\end{minipage} & \begin{minipage}[b]{\linewidth}\raggedright
Phoneme
\end{minipage} & \begin{minipage}[b]{\linewidth}\raggedright
Application Rule
\end{minipage} & \begin{minipage}[b]{\linewidth}\raggedright
Example
\end{minipage} \\
\midrule\noalign{}
\endhead
\bottomrule\noalign{}
\endlastfoot
\textbf{c} & /k/ & Soft, unaspirated velar plosive. & cat, cum, cup,
bac \\
\textbf{k} & /kʰ/ & Hard, aspirated velar plosive. & king, kiip, kic \\
\textbf{ci} / \textbf{ce} & /tʃ/ & Voiceless palato-alveolar affricate.
The vowel functions as an orthographic trigger. & ciald, cier, ceic
(check) \\
\textbf{u} & /w/ & Semivowel substitution for \emph{w}. & uoter, uid \\
\textbf{i} & /j/ & Semivowel substitution for \emph{y}. & ies, iir \\
\textbf{t} & /t/ or /θ/ & Alveolar plosive; assumes the voiceless dental
fricative locally. & tinc, trii \\
\textbf{d} & /d/ or /ð/ & Alveolar plosive; assumes the voiced dental
fricative locally. & de, dat \\
\end{longtable}
}

\begin{quote}
\textbf{Aspiration Differentiation (\texttt{k} vs \texttt{c}):} The
aspirated /kʰ/ is strictly represented by \texttt{k}, whereas the
unaspirated variant uses \texttt{c}. Following sibilants (e.g., /s/),
where aspiration naturally diminishes, \texttt{c} is mandated (e.g.,
\emph{scan}, \emph{scip}).
\end{quote}

\subsubsection{\texorpdfstring{2.3.1 Lexical Prosody Marker
(\texttt{\textquotesingle{}})}{2.3.1 Lexical Prosody Marker (\textquotesingle)}}\label{lexical-prosody-marker}

Because MinEnglish eliminates the redundant double consonants and silent
terminal vowels often used as prosodic cues in natural English, lexical
stress can become ambiguous for computational parsers. The apostrophe
(\texttt{\textquotesingle{}}) is instituted as an optional prosody
marker, to be affixed immediately preceding the stressed vowel in
pedagogy or lexicography.

\begin{itemize}
\tightlist
\item
  \texttt{ecsp\textquotesingle{}ensiv} (Stress alignment on the initial
  \emph{e})
\item
  \texttt{iuniv\textquotesingle{}ersiti} (Stress alignment on the medial
  \emph{e})
\item
  \texttt{compi\textquotesingle{}uuter} (Stress alignment on the first
  \emph{u})
\end{itemize}

\subsubsection{2.4 Digraphs}\label{digraphs}

{\def\LTcaptype{} % do not increment counter
\begin{longtable}[]{@{}lll@{}}
\toprule\noalign{}
Digraph & Phoneme & Example \\
\midrule\noalign{}
\endhead
\bottomrule\noalign{}
\endlastfoot
\textbf{sh} & /ʃ/ & ship → ship \\
\textbf{ng} & /ŋ/ & song → song \\
\textbf{zh} & /ʒ/ & vision → vizhun \\
\end{longtable}
}

\begin{quote}
Note: The standard English \emph{ch} (/tʃ/) is not utilized. It is
strictly encoded as \texttt{ci} to maintain morphemic consistency.
\end{quote}

\subsubsection{2.5 Orthographic Normalization
Constraints}\label{orthographic-normalization-constraints}

To enforce the phonetic isomorphism principle, historical orthographic
artifacts must be algorithmically stripped.

{\def\LTcaptype{} % do not increment counter
\begin{longtable}[]{@{}
  >{\raggedright\arraybackslash}p{(\linewidth - 4\tabcolsep) * \real{0.3579}}
  >{\raggedright\arraybackslash}p{(\linewidth - 4\tabcolsep) * \real{0.4211}}
  >{\raggedright\arraybackslash}p{(\linewidth - 4\tabcolsep) * \real{0.2211}}@{}}
\toprule\noalign{}
\begin{minipage}[b]{\linewidth}\raggedright
Historical Pattern
\end{minipage} & \begin{minipage}[b]{\linewidth}\raggedright
MinEnglish Normalization
\end{minipage} & \begin{minipage}[b]{\linewidth}\raggedright
Example
\end{minipage} \\
\midrule\noalign{}
\endhead
\bottomrule\noalign{}
\endlastfoot
Silent \textbf{e}, \textbf{k}, \textbf{g}, \textbf{gh} & Nullified &
meic, nou, noum, nait \\
\textbf{gh} or \textbf{ph} rendered as /f/ & Regressive substitution to
\textbf{f} & laf, foun \\
\textbf{ck} cluster & Resolved to \textbf{c} or \textbf{k} by aspiration
& bac, kic \\
\textbf{x} cluster & Decomposed to \textbf{cs} & bocs, necst \\
\textbf{-tion} / \textbf{-sion} & Phoneticized to \textbf{-shun} /
\textbf{-zhun} & neishun, vizhun \\
\textbf{-tch} & Phoneticized to \textbf{-ci} & maci, caci \\
Double consonants & Replaced with unitary equivalents & leter, hapi \\
Soft \textbf{c} (/s/) & Regressive substitution to \textbf{s} & sel,
siti \\
\end{longtable}
}

\begin{center}\rule{0.5\linewidth}{0.5pt}\end{center}

\newpage

\subsection{3. Morphology and Syntax}\label{morphology-and-syntax}

\begin{quote}
\textbf{Core Constraint: Morphological Invariance.} Lexical items never
undergo morphological derivation or inflection. Grammatical
relationships are managed exclusively through standardized operators and
rigid syntactic ordering.
\end{quote}

\subsubsection{3.1 Nouns and Quantitative
Prefixes}\label{nouns-and-quantitative-prefixes}

To eliminate the biological error-amplifier of plural suffixes (where
the loss of a terminal sibilant in a noisy channel alters the
mathematical quantity of a sentence), MinEnglish mandates explicit
numerical prefixing. \textbf{Morphological plurals are abolished.}

{\def\LTcaptype{} % do not increment counter
\begin{longtable}[]{@{}lll@{}}
\toprule\noalign{}
Prefix & Semantic Value & Example \\
\midrule\noalign{}
\endhead
\bottomrule\noalign{}
\endlastfoot
\texttt{1} & Unitary / Singular & \texttt{1cat} (a cat) \\
\texttt{3} & Exact enumeration & \texttt{3cat} (three cats) \\
\texttt{*} & Universal / Generic set & \texttt{*cat} (cats, across all
domains) \\
\texttt{\textasciitilde{}5} & Approximate enumeration &
\texttt{\textasciitilde{}5cat} (approximately five cats) \\
\texttt{0} & Null set & \texttt{0cat} (zero cats) \\
\end{longtable}
}

\begin{quote}
\textbf{Absence of Articles:} The quantitative operator \texttt{1}
functionally replaces the indefinite article. Definiteness is
established contextually or via demonstrative pronouns.
\end{quote}

\subsubsection{3.1.1 Numerical Processing
Efficiency}\label{numerical-processing-efficiency}

To minimize phonological bandwidth, integers exceeding 99 bypass
categorical naming structures and are processed as pure digit strings.

\begin{itemize}
\tightlist
\item
  \texttt{3456} = \texttt{trii\ for\ faiv\ sics}
\item
  \texttt{2026} = \texttt{tuu\ ziro\ tuu\ sics}
\end{itemize}

\subsubsection{3.1.2 Proper Noun
Encapsulation}\label{proper-noun-encapsulation}

Phonetic normalization of global proper nouns creates destructive
variations that impair search retrieval and semantic persistence.

\textbf{Rule:} Proper Nouns (specific entities, geographical locations,
nominal brands) retain their canonical native orthography. First-letter
capitalization acts as an escape character, signaling to the parser that
local phonetic rules are suspended for the duration of the token.

\begin{quote}
\textbf{Strict Lowercase Axiom:} Because capitalization is an active
mathematical operator (the Escape key), MinEnglish strictly prohibits
capitalizing the first letter of sentences. Capitalizing a
sentence-starting word would force parsers to guess whether the word is
an escaped proper noun or standard orthographic decoration, destroying
mathematical purity. All standard MinEnglish text is strictly lowercase.
\end{quote}

\begin{itemize}
\tightlist
\item
  \texttt{person\ Isabella} (Denotes the individual Isabella;
  phonetically preserved).
\item
  \texttt{cuntri\ France} (Denotes the sovereign state of France).
\end{itemize}

\subsubsection{3.1.3 Fractional and Decimal
Operators}\label{fractional-and-decimal-operators}

Spoken mathematics is standardized via the \texttt{/} (slais) and
\texttt{.} (point) operators, providing a syntactically uniform method
for subdividing physical objects or abstract metrics.

\begin{itemize}
\tightlist
\item
  \textbf{Fractions:} \texttt{1/2} (uan slais tuu) = One-half.
  \texttt{3/4apel} (trii slais for apel) = Three-quarters of an apple.
\item
  \textbf{Decimals:} \texttt{0.5} (ziro point faiv). \texttt{99.9} (nain
  nain point nain).
\end{itemize}

\subsubsection{3.1.4 Token Optionality
Rules}\label{token-optionality-rules}

Three categories of token may be omitted when their value is recoverable
from context. \textbf{Parsers must implement the following defaults when
a token is absent:}

\paragraph{\texorpdfstring{Rule A --- Implicit Singular (\texttt{1}
optional)}{Rule A --- Implicit Singular (1 optional)}}\label{rule-a-implicit-singular-1-optional}

When no quantitative prefix is present before a noun, the parser
defaults to singular (\texttt{1}). The explicit \texttt{1} is required
only when disambiguation from a preceding plural context is necessary.

{\def\LTcaptype{} % do not increment counter
\begin{longtable}[]{@{}lll@{}}
\toprule\noalign{}
Explicit & Implicit & Meaning \\
\midrule\noalign{}
\endhead
\bottomrule\noalign{}
\endlastfoot
\texttt{1cat\ sit\ on\ 1mat} & \texttt{cat\ sit\ on\ mat} & A cat sits
on a mat \\
\texttt{1dog\ run\ in\ 1parc} & \texttt{dog\ run\ in\ parc} & A dog runs
in a park \\
\end{longtable}
}

\begin{quote}
\textbf{Note:} If the previous noun was plural, explicit \texttt{1} is
required to re-establish singular: \texttt{3cat\ run,\ den\ 1cat\ stop.}
\end{quote}

\paragraph{\texorpdfstring{Rule B --- Implicit First Person (\texttt{i}
optional)}{Rule B --- Implicit First Person (i optional)}}\label{rule-b-implicit-first-person-i-optional}

When the subject is first-person singular and the sentence structure is
declarative, the pronoun \texttt{i} may be omitted. The parser defaults
to \texttt{i} when no subject precedes the verb block.

{\def\LTcaptype{} % do not increment counter
\begin{longtable}[]{@{}
  >{\raggedright\arraybackslash}p{(\linewidth - 4\tabcolsep) * \real{0.2812}}
  >{\raggedright\arraybackslash}p{(\linewidth - 4\tabcolsep) * \real{0.2656}}
  >{\raggedright\arraybackslash}p{(\linewidth - 4\tabcolsep) * \real{0.4531}}@{}}
\toprule\noalign{}
\begin{minipage}[b]{\linewidth}\raggedright
Explicit
\end{minipage} & \begin{minipage}[b]{\linewidth}\raggedright
Implicit
\end{minipage} & \begin{minipage}[b]{\linewidth}\raggedright
Meaning
\end{minipage} \\
\midrule\noalign{}
\endhead
\bottomrule\noalign{}
\endlastfoot
\texttt{i1laic\ cat} & \texttt{laic\ cat} & I like a cat \\
\texttt{i-1d\ gou\ tu\ stor} & \texttt{-1d\ gou\ tu\ stor} & I went to
the store yesterday \\
\end{longtable}
}

\begin{quote}
\textbf{Disambiguation:} If a third-person pronoun or noun immediately
precedes the verb, \texttt{i} is not assumed. The implicit-I rule only
fires when the verb begins the sentence.
\end{quote}

\paragraph{Rule C --- Implicit Present Tense (temporal marker
optional)}\label{rule-c-implicit-present-tense-temporal-marker-optional}

A verb with no temporal prefix is already defined as present continuous.
Additionally, the discrete present marker \texttt{1} (as in
\texttt{h1laic}) is also optional when the subject is explicit and no
tense ambiguity exists. The \texttt{1} attachment is required only for
disambiguation or emphasis.

{\def\LTcaptype{} % do not increment counter
\begin{longtable}[]{@{}lll@{}}
\toprule\noalign{}
Explicit & Implicit & Meaning \\
\midrule\noalign{}
\endhead
\bottomrule\noalign{}
\endlastfoot
\texttt{h1laic\ 1cat} & \texttt{h\ laic\ cat} & He likes a cat \\
\texttt{i1run\ fast} & \texttt{run\ fast} & I run fast \\
\end{longtable}
}

\subsubsection{3.2 Verbs --- Temporal
Mechanics}\label{verbs-temporal-mechanics}

Irregular conjugations are functionally inefficient. MinEnglish
regulates time via temporal operators that target an absolute
mathematical timeline. A verb without a temporal operator defaults to
present continuous.

{\def\LTcaptype{} % do not increment counter
\begin{longtable}[]{@{}
  >{\raggedright\arraybackslash}p{(\linewidth - 4\tabcolsep) * \real{0.1200}}
  >{\raggedright\arraybackslash}p{(\linewidth - 4\tabcolsep) * \real{0.4100}}
  >{\raggedright\arraybackslash}p{(\linewidth - 4\tabcolsep) * \real{0.4700}}@{}}
\toprule\noalign{}
\begin{minipage}[b]{\linewidth}\raggedright
Operator
\end{minipage} & \begin{minipage}[b]{\linewidth}\raggedright
Temporal Value
\end{minipage} & \begin{minipage}[b]{\linewidth}\raggedright
Example
\end{minipage} \\
\midrule\noalign{}
\endhead
\bottomrule\noalign{}
\endlastfoot
\emph{(null)} & Present continuous (default) & \texttt{iit} (eating) \\
\texttt{1} & Present discrete (attachment operator) & \texttt{i1iit} (I
eat) \\
\texttt{*} & Habitual / Universal continuity & \texttt{*iit} (eats
perpetually) \\
\texttt{-0s} & Immediate past vector (zero offset) & \texttt{-0s\ sii}
(just perceived) \\
\texttt{-1d} & Past vector (one day offset) & \texttt{-1d\ iit} (ate
yesterday) \\
\texttt{:5Y} & Duration scalar (five years) & \texttt{:5Y\ studi} (has
studied for five years) \\
\texttt{1d} & Future vector (\texttt{+} accepted but optional) &
\texttt{1d\ gou} (will go tomorrow) \\
\texttt{2026-12-25} & Absolute timestamp & \texttt{2026-12-25\ gou}
(transit scheduled for Dec 25) \\
\end{longtable}
}

\begin{quote}
\textbf{Direction Convention:} The \texttt{-} prefix marks past vectors
(required). Future offsets default to positive --- a bare numeric offset
(\texttt{1d}, \texttt{2h}, \texttt{1Y}) always points forward. The
\texttt{:} operator spans a duration rather than pointing to a discrete
moment.
\end{quote}

\subsubsection{3.2.1 Temporal Stacking (Complex
Timelines)}\label{temporal-stacking-complex-timelines}

Natural language manages chronological depth via ``perfect'' tenses,
increasing syntactic distance between subject and action. MinEnglish
resolves deep time by permitting left-to-right stacking of temporal
operators, enabling precise algebraic timeline construction.

\begin{itemize}
\tightlist
\item
  \texttt{-0s\ -1d\ iit} = I just-now yesterday ate (I had already eaten
  by yesterday's timeline).
\item
  \texttt{-1d\ 2h\ iit} = Yesterday, offset by +2 hours, I ate.
\end{itemize}

\subsubsection{3.3 Modifiers and Scalar
Intensity}\label{modifiers-and-scalar-intensity}

\begin{itemize}
\tightlist
\item
  \textbf{Adjectival Ordering:} Adjectives strictly post-modify nouns
  (\texttt{cat\ big}).
\item
  \textbf{Adverbial Ordering:} Adverbs post-modify verbs
  (\texttt{run\ fast}).
\item
  \textbf{Comparative Mechanics:} Driven by scalar progression
  \texttt{mor} and \texttt{most} (\texttt{cat\ mor\ big}).
\item
  \textbf{Intensity Combinators (\texttt{\textgreater{}},
  \texttt{\textgreater{}\textgreater{}}, \texttt{\textless{}},
  \texttt{\textless{}\textless{}}):} Adjectives and adverbs lack
  morphological escalation (no ``-er'' or ``-est''). Instead, magnitude
  is modulated using absolute mathematical prefixes that attach directly
  to the front of their target.

  \begin{itemize}
  \tightlist
  \item
    \texttt{\textgreater{}big} = Positive magnification (very big)
  \item
    \texttt{\textgreater{}\textgreater{}big} = Extreme maximization
    (massive, huge)
  \item
    \texttt{\textless{}big} = Negative attenuation (slightly big)
  \item
    \texttt{\textless{}\textless{}big} = Extreme minimization (tiny,
    very small)
  \end{itemize}
\end{itemize}

\subsubsection{\texorpdfstring{3.3.1 Universal Predicators (\texttt{du}
/
\texttt{meic})}{3.3.1 Universal Predicators (du / meic)}}\label{universal-predicators-du-meic}

Because MinEnglish permits any token to accept an operational prefix,
inappropriate syntactic mapping can yield semantic anomalies. The
universal predicators \texttt{du} (experiential interaction) and
\texttt{meic} (generative action) serve as formal verbalizers for
noun-class tokens.

\begin{itemize}
\tightlist
\item
  \texttt{du\ parc} = I interact with the park (replaces the ambiguous
  \texttt{parc}).
\item
  \texttt{s-1d\ meic\ fuud} = She generated food (replaces the anomalous
  \texttt{s-1d\ fuud}).
\end{itemize}

\subsubsection{3.4 Pronominal Indexing}\label{pronominal-indexing}

Pronouns operate as highly compressed, single-character spatial indices.
Plurality logic inherits the identical numerical prefixing system
utilized by canonical nouns.

{\def\LTcaptype{} % do not increment counter
\begin{longtable}[]{@{}ll@{}}
\toprule\noalign{}
MinEnglish Singular & MinEnglish Plural \\
\midrule\noalign{}
\endhead
\bottomrule\noalign{}
\endlastfoot
\textbf{i} (I) & \texttt{5i} (We five) \\
\textbf{u} (You) & \texttt{3u} (You three) \\
\textbf{h} (He) & \texttt{4h} (They four males) \\
\textbf{s} (She) & \texttt{21s} (They 21 females) \\
\textbf{t} (It) & \texttt{*t} (They things) \\
\end{longtable}
}

\begin{quote}
\textbf{Genitive Case Marking:} Possession is indicated universally via
the \texttt{\textquotesingle{}s} suffix with no exceptions. (e.g.,
\texttt{i\textquotesingle{}s}, \texttt{u\textquotesingle{}s},
\texttt{h\textquotesingle{}s}, \texttt{s\textquotesingle{}s},
\texttt{t\textquotesingle{}s}).
\end{quote}

\subsubsection{3.4.1 Syntactic Agglutination
(Attachment)}\label{syntactic-agglutination-attachment}

For maximum data throughput, the single-character pronominal index
agglutinates physically to the succeeding temporal operator, eliminating
whitespace parsing overhead.

\begin{Shaded}
\begin{Highlighting}[]
\NormalTok{h1laic    (he likes)}
\NormalTok{s1run     (she runs)}
\NormalTok{u:2h uorc (you worked spanning two hours)}
\NormalTok{s*run     (she runs generically)}
\end{Highlighting}
\end{Shaded}

\subsubsection{3.4.2 Auditory Disambiguation (Prefix
Positioning)}\label{auditory-disambiguation-prefix-positioning}

Agglutination can generate overlapping vowel harmonics (e.g.,
\texttt{in} or \texttt{if\ i*no}), increasing transmission error rates.
To preserve auditory clarity boundaries, pronominal indices have
flexible positioning relative to the temporal operator.

\begin{itemize}
\tightlist
\item
  Canonical: \texttt{i-1d} (I yesterday)
\item
  Permuted: \texttt{1d-i} (Yesterday I; mitigates harmonic overlap after
  preceding vowels)
\end{itemize}

\subsubsection{3.4.3 Pronominal Deference}\label{pronominal-deference}

By default, the syntactic structure is declarative and absolute. To
inject social deference or epistemological uncertainty, the
approximation operator \texttt{\textasciitilde{}} is bound to the
pronoun.

\begin{itemize}
\tightlist
\item
  \texttt{u\ giv} = Absolute directive (You give).
\item
  \texttt{\textasciitilde{}u\ giv} = Deferred directive (Could you
  please give).
\item
  \texttt{\textasciitilde{}tinc} = Epistemological deferment (It seems
  to me).
\end{itemize}

\subsubsection{3.5 Negation}\label{negation}

Negation is uniquely expressed by attaching the \texttt{!} prefix
directly to the target word (typically the verb or adjective) without a
space.

\begin{itemize}
\tightlist
\item
  \texttt{!laic\ cat} = I do not like the cat.
\item
  \texttt{u1d\ !cum} = You will not come tomorrow.
\item
  \texttt{*h\ !can\ spiic} = None of them can speak.
\item
  \texttt{bi\ !gud} = It is not good.
\end{itemize}

\subsubsection{3.6 Sentence Structure \& Conjunction
Reduction}\label{sentence-structure-conjunction-reduction}

\begin{quote}
\textbf{Subject + Time-Verb + Number-Object + Modifiers}
\end{quote}

\begin{Shaded}
\begin{Highlighting}[]
\NormalTok{laic cat \textgreater{}big.            → I like a very big cat.}
\NormalTok{*h{-}1d bai 3buc.              → They bought 3 books yesterday.}
\end{Highlighting}
\end{Shaded}

\subsubsection{3.6.1 Conjunction Reduction}\label{conjunction-reduction}

If multiple actions share the same subject, you can drop the subject
after \texttt{an} or \texttt{or} (Conjunction Reduction):

\begin{itemize}
\tightlist
\item
  \texttt{iit\ an\ drinc} = I eat and (I) drink.
\item
  \texttt{iit\ an\ h1drinc} = I eat and he drinks.
\end{itemize}

\subsubsection{3.7 Modality (can, must,
should)}\label{modality-can-must-should}

Modality verbs follow the time prefix, creating a composite verb phrase:

{\def\LTcaptype{} % do not increment counter
\begin{longtable}[]{@{}
  >{\raggedright\arraybackslash}p{(\linewidth - 4\tabcolsep) * \real{0.1159}}
  >{\raggedright\arraybackslash}p{(\linewidth - 4\tabcolsep) * \real{0.3188}}
  >{\raggedright\arraybackslash}p{(\linewidth - 4\tabcolsep) * \real{0.5652}}@{}}
\toprule\noalign{}
\begin{minipage}[b]{\linewidth}\raggedright
Modal
\end{minipage} & \begin{minipage}[b]{\linewidth}\raggedright
Meaning
\end{minipage} & \begin{minipage}[b]{\linewidth}\raggedright
Example
\end{minipage} \\
\midrule\noalign{}
\endhead
\bottomrule\noalign{}
\endlastfoot
\textbf{can} & ability / permission & \texttt{can\ iit} = I can eat \\
\textbf{must} & necessity / obligation & \texttt{u1d\ must\ cum} = You
must come tomorrow \\
\textbf{shud} & advice / should & \texttt{*h\ shud\ lern} = They should
learn \\
\textbf{\textasciitilde can} & probability / might &
\texttt{i1d\ \textasciitilde{}can\ gou} = I might go tomorrow \\
\end{longtable}
}

\subsubsection{3.8 Passive Voice}\label{passive-voice}

Formed by moving the Object to the front, and following the verb with
\texttt{bai} (by).

{\def\LTcaptype{} % do not increment counter
\begin{longtable}[]{@{}
  >{\raggedright\arraybackslash}p{(\linewidth - 4\tabcolsep) * \real{0.1196}}
  >{\raggedright\arraybackslash}p{(\linewidth - 4\tabcolsep) * \real{0.3261}}
  >{\raggedright\arraybackslash}p{(\linewidth - 4\tabcolsep) * \real{0.5543}}@{}}
\toprule\noalign{}
\begin{minipage}[b]{\linewidth}\raggedright
Voice
\end{minipage} & \begin{minipage}[b]{\linewidth}\raggedright
MinEnglish Structure
\end{minipage} & \begin{minipage}[b]{\linewidth}\raggedright
Example
\end{minipage} \\
\midrule\noalign{}
\endhead
\bottomrule\noalign{}
\endlastfoot
\textbf{Active} & Subj + Time-Verb + Obj & \texttt{boi\ kic\ bol} (A boy
just kicked a ball) \\
\textbf{Passive} & Obj + Time-Verb + \texttt{bai} + Subj &
\texttt{bol\ kic\ bai\ boi} (A ball was just kicked by a boy) \\
\end{longtable}
}

\subsubsection{\texorpdfstring{3.8.1 The Direct Object Marker
(\texttt{tu})}{3.8.1 The Direct Object Marker (tu)}}\label{the-direct-object-marker-tu}

Because MinEnglish removes noun-verb class distinctions (any word can
technically be a verb with \texttt{1}), complex sentences can blur the
line between a chained verb and a direct object noun. \textbf{Rule:}
When a noun could be mistaken for a verb, prefix it with the transitive
operator \texttt{tu} (acting like Esperanto's -n accusative).

\begin{itemize}
\tightlist
\item
  \texttt{laic\ tu\ cat} = I like the cat (Prevents reading as ``I like
  and I cat'').
\item
  \texttt{s1giv\ tu\ h\ buc} = She gives him a book.
\end{itemize}

\subsubsection{3.9 Connectors, Prepositions, Tone
Markers}\label{connectors-prepositions-tone-markers}

\begin{itemize}
\tightlist
\item
  \textbf{Logic:} \texttt{an} / \texttt{\&} (and), \texttt{or} /
  \texttt{\textbar{}} (or), \texttt{\^{}} (xor), \texttt{but} (but),
  \texttt{if} (if), \texttt{cos} (because), \texttt{sou} (so),
  \texttt{den} (then).
\item
  \textbf{Shortcuts:} \texttt{@} (at), \texttt{\textgreater{}} (over /
  more than), \texttt{\textless{}} (under / less than).
\item
  \textbf{Tone:} \texttt{!} (emphasis), \texttt{?} (question),
  \texttt{..} (hesitation), \texttt{\^{}} (sarcasm).
\end{itemize}

\begin{quote}
\textbf{Symbolic Connectors:} \texttt{\&} and \texttt{\textbar{}} are
fully equivalent to \texttt{an} and \texttt{or} respectively, and are
preferred in dense or technical text. In speech, both forms are
pronounced identically (\texttt{an}, \texttt{or}).
\end{quote}

\subsubsection{\texorpdfstring{3.9.1 Terminal Punctuation \& Negation
(\texttt{!},
\texttt{?})}{3.9.1 Terminal Punctuation \& Negation (!, ?)}}\label{terminal-punctuation-negation}

To maintain familiarity and avoid syntactic overcomplication, MinEnglish
uses the standard English terminal punctuation marks (\texttt{?} and
\texttt{!}) located at the end of the sentence clause.

However, the \texttt{!} symbol has a dual function depending on
whitespace spacing:

\begin{itemize}
\item
  \textbf{Exclamation (Space/End):} When placed at the end of a sentence
  or separated by a space, it acts natively as an exclamation. Doubling
  the terminal marks (\texttt{!!} or \texttt{??}) is supported and
  explicitly defined as an aggressive, warning, or obnoxious tone
  multiplier.
\item
  \textbf{Negation Prefix (Attached):} When attached directly to the
  front of a word with \emph{no space}, it acts as the logical NOT
  operator.
\item
  \texttt{uen\ u1gou??} = \textbf{When} do you go?! (Obnoxious
  Interrogation)
\item
  \texttt{uoci!!} = Watch out! (Aggressive Warning)
\item
  \texttt{i\ !laic\ cat.} = I do \textbf{not} like cats. (Negation)
\item
  \texttt{!gud} = Not good.
\item
  \texttt{huu}, \texttt{uat}, \texttt{uai}, \texttt{haum} (who, what,
  why, how much-many).
\end{itemize}

\subsubsection{\texorpdfstring{3.9.2 Clause Anchors (\texttt{\{} /
\texttt{\}})}{3.9.2 Clause Anchors (\{ / \})}}\label{clause-anchors}

Because MinEnglish removes grammatical cases and subordinate conjunction
fluff, deeply center-embedded sentences will physically overwhelm human
short-term memory (the phonological loop). \textbf{Rule:} When nesting 2
or more clauses, the speaker \emph{must} use the explicit mathematical
scope operators \texttt{\{} and \texttt{\}} to define the boundaries of
the nested subset.

\begin{itemize}
\tightlist
\item
  \texttt{man\ \{\ dat\ dog\ \{\ dat\ ciald\ luv\ \}\ -1d\ bait\ \}\ run\ auei.}
\item
  \emph{Translation:} The man {[}who the dog {[}that the child loves{]}
  bit{]} ran away.
\end{itemize}

\begin{quote}
\textbf{Phonetic Pronunciation (\texttt{co} / \texttt{oc}):} In purely
auditory communication where typographical braces cannot be vocalized
efficiently, the strings \texttt{co} (clause open) and \texttt{oc}
(clause close) are the mandatory phonetic equivalents to \texttt{\{} and
\texttt{\}}.

\begin{itemize}
\tightlist
\item
  \textbf{Auditory Mapping:}
  \texttt{man\ co\ dat\ dog\ co\ dat\ ciald\ luv\ oc\ -1d\ bait\ oc\ run\ auei.}
\end{itemize}
\end{quote}

\subsubsection{3.9.3 Arithmetic \& Logical Operator
Inventory}\label{arithmetic-logical-operator-inventory}

MinEnglish formally adopts five symbolic operators as first-class
syntactic tokens. \texttt{/} is reserved for fractions and math (see
3.1.4) and is \textbf{not} a compound word operator.

{\def\LTcaptype{} % do not increment counter
\begin{longtable}[]{@{}
  >{\raggedright\arraybackslash}p{(\linewidth - 6\tabcolsep) * \real{0.0930}}
  >{\raggedright\arraybackslash}p{(\linewidth - 6\tabcolsep) * \real{0.1047}}
  >{\raggedright\arraybackslash}p{(\linewidth - 6\tabcolsep) * \real{0.4884}}
  >{\raggedright\arraybackslash}p{(\linewidth - 6\tabcolsep) * \real{0.3140}}@{}}
\toprule\noalign{}
\begin{minipage}[b]{\linewidth}\raggedright
Operator
\end{minipage} & \begin{minipage}[b]{\linewidth}\raggedright
Name
\end{minipage} & \begin{minipage}[b]{\linewidth}\raggedright
Primary Role
\end{minipage} & \begin{minipage}[b]{\linewidth}\raggedright
Example
\end{minipage} \\
\midrule\noalign{}
\endhead
\bottomrule\noalign{}
\endlastfoot
\texttt{+} & Plus & Future offset (optional; bare N = default) &
\texttt{1d\ gou} = will go tomorrow \\
\texttt{-} & Minus & Past offset (required) & \texttt{-1d\ iit} = ate
yesterday \\
\texttt{*} & Star & Universal / generic quantifier & \texttt{*cat} = all
cats \\
\texttt{\&} & Ampersand & Logical AND (≡ \texttt{an}) &
\texttt{laic\ cat\ \&\ dog} \\
\texttt{\textbar{}} & Pipe & Logical OR (≡ \texttt{or}) & \\
\texttt{\^{}} & Caret & Logical XOR (Exclusive OR) &
\texttt{iit\ cat\ \^{}\ dog} (one only) \\
\end{longtable}
}

\begin{quote}
\texttt{*i} means ``all of us'' (universal quantifier on first-person
pronoun) --- not habitual verb. When \texttt{*} prefixes a pronoun it
quantifies it; when it prefixes a verb it marks habitual action. Parsers
distinguish by token class.
\end{quote}

\subsubsection{3.9.4 Preposition Verbs}\label{preposition-verbs}

To avoid repeating the verb ``to be'' (\texttt{bi}), prepositions can
act as verbs simply by attaching a time prefix to them.

\begin{itemize}
\tightlist
\item
  \texttt{haus\ in\ siti} = The house is in a city (\texttt{in} acts as
  the verb ``to be in'').
\item
  \texttt{cat\ -1d\ on\ mat} = The cat was on a mat yesterday.
\end{itemize}

\subsubsection{3.10 Abjad Mode (Extreme
Shorthand)}\label{abjad-mode-extreme-shorthand}

For ultra-fast texting or constrained bandwidth, MinEnglish formally
defines ``Abjad Mode''. \textbf{Rule:} Drop all vowels \emph{unless} the
word starts with a vowel or doing so breaks a core operator.

\begin{itemize}
\tightlist
\item
  Standard / English: \texttt{tinc\ u\ bi\ \textgreater{}liit} (I think
  you are very late)
\item
  Abjad Mode: \texttt{tnc\ u\ b\ \textgreater{}lt}
\end{itemize}

\subsubsection{3.11 Semantic Idiom
Standardization}\label{semantic-idiom-standardization}

English relies on illogical phrasal verbs (``Look out!'' = beware;
``Give up'' = surrender). MinEnglish rejects phrasal idioms entirely. A
direct translation of ``Look out'' (\texttt{luc\ aut}) is strictly
literal: point your eyes outside. \textbf{Rule:} Idioms must be
translated to their literal action equivalent.

\begin{itemize}
\tightlist
\item
  ``Look out!'' → \texttt{uoci!!} (Watch intensely!)
\item
  ``Give up'' → \texttt{stop\ tu\ trai} (Stop to try)
\end{itemize}

\subsubsection{3.12 Orthographic Stress
Engine}\label{orthographic-stress-engine}

To solve the ambiguity of spoken MinEnglish without relying on the
optional \texttt{\textquotesingle{}} marker, standardizes default
pronunciation logic for parsers and speakers. \textbf{Rule:} Always
stress the \textbf{first vowel} of a root word \emph{unless} it is a
recognized prefix-suffix. In compound words, stress the first syllable
of the primary noun.

\begin{itemize}
\tightlist
\item
  \texttt{c\textquotesingle{}ompiuuter} (Stress the O)
\item
  \texttt{b\textquotesingle{}ioloji} (Stress the I)
\end{itemize}

\subsubsection{3.13 Echo-Tagging (Aviation/Noisy
Protocol)}\label{echo-tagging-aviationnoisy-protocol}

Claude Shannon's Information Theory proves that ``noise'' requires
``redundancy'' to fix. MinEnglish's maximum data density strips all
redundancy (\texttt{3boi\ uoc} marks plural only once). In high-noise
environments (construction, radios, wind), MinEnglish fails
catastrophically if one syllable drops. \textbf{Rule:} For critical
transmissions, a speaker may artificially inflate Shannon redundancy to
100\% by ``echoing'' the prefix as a full word \emph{after} the root.

\begin{itemize}
\tightlist
\item
  Standard: \texttt{3boi\ uoc} (Three boys are walking)
\item
  Echo-Tagged: \texttt{3boi\ trii\ uoc\ nau} (Three boys three
  are-walking now)
\end{itemize}

\subsubsection{3.14 Formal Syntax (Backus-Naur
Form)}\label{formal-syntax-backus-naur-form}

The absolute regularity of MinEnglish allows it to be algorithmically
validated. Below is the strict parsing configuration for a standard
declarative sentence.

\begin{Shaded}
\begin{Highlighting}[]
\NormalTok{\textless{}sentence\textgreater{}      ::= [\textless{}subject\_block\textgreater{}] \textless{}verb\_block\textgreater{} [\textless{}object\_block\textgreater{}] [\textless{}modifier\_block\textgreater{}]...}
\NormalTok{\textless{}subject\_block\textgreater{} ::= \textless{}noun\_phrase\textgreater{} | \textless{}pronoun\_phrase\textgreater{}}
\NormalTok{\textless{}object\_block\textgreater{}  ::= \textless{}noun\_phrase\textgreater{} | \textless{}pronoun\_phrase\textgreater{}}

\NormalTok{\textless{}noun\_phrase\textgreater{}   ::= \textless{}quantity\textgreater{} \textless{}noun\_root\textgreater{} [\textless{}adjective\_phrase\textgreater{}]}
\NormalTok{\textless{}pronoun\_phrase\textgreater{}::= [\textless{}quantity\textgreater{}] \textless{}pronoun\_root\textgreater{}}
\NormalTok{\textless{}adjective\_phrase\textgreater{} ::= [\textless{}intensity\textgreater{}] \textless{}adjective\_root\textgreater{}}

\NormalTok{\textless{}verb\_block\textgreater{}    ::= [\textless{}temporal\_vector\textgreater{}] \textless{}verb\_root\textgreater{} [\textless{}adverb\_phrase\textgreater{}]}
\NormalTok{\textless{}temporal\_vector\textgreater{} ::= \textless{}time\_prefix\textgreater{} | \textless{}duration\_prefix\textgreater{} | \textless{}unix\_timestamp\textgreater{}}

\NormalTok{\textless{}clause\textgreater{}        ::= "\{" \textless{}connective\textgreater{} \textless{}sentence\textgreater{} "\}"}
\end{Highlighting}
\end{Shaded}

\begin{center}\rule{0.5\linewidth}{0.5pt}\end{center}

\newpage

\subsection{4. Base Lexicon}\label{base-lexicon}

\subsubsection{4.1 Fundamentals}\label{fundamentals}

{\def\LTcaptype{} % do not increment counter
\begin{longtable}[]{@{}ll@{}}
\toprule\noalign{}
English & MinEnglish \\
\midrule\noalign{}
\endhead
\bottomrule\noalign{}
\endlastfoot
be / is & \textbf{bi} \\
do & \textbf{du} \\
come & \textbf{cum} \\
take & \textbf{teic} \\
see & \textbf{sii} \\
think & \textbf{tinc} \\
want / need & \textbf{uant} / \textbf{niid} \\
have & \textbf{hav} \\
go & \textbf{gou} \\
make & \textbf{meic} \\
give & \textbf{giv} \\
know & \textbf{nou} \\
say / tell & \textbf{sei} / \textbf{tel} \\
use / find & \textbf{iuz} / \textbf{faind} \\
\end{longtable}
}

\subsubsection{4.2 Time, Numbers, Nature}\label{time-numbers-nature}

{\def\LTcaptype{} % do not increment counter
\begin{longtable}[]{@{}ll@{}}
\toprule\noalign{}
English & MinEnglish \\
\midrule\noalign{}
\endhead
\bottomrule\noalign{}
\endlastfoot
1, 2, 3, 4, 5 & uan, tuu, trii, for, faiv \\
6, 7, 8, 9, 0 & sics, seven, eit, nain, ziro \\
10, 100, 1k & ten, hundrid, tauzund \\
time & \textbf{taim} \\
day / night & \textbf{dei} / \textbf{nait} \\
month / year & \textbf{munt} / \textbf{iir} \\
sun / moon & \textbf{sun} / \textbf{muun} \\
rain / snow & \textbf{rein} / \textbf{snou} \\
\end{longtable}
}

\subsubsection{4.3 People \& Body}\label{people-body}

{\def\LTcaptype{} % do not increment counter
\begin{longtable}[]{@{}ll@{}}
\toprule\noalign{}
English & MinEnglish \\
\midrule\noalign{}
\endhead
\bottomrule\noalign{}
\endlastfoot
person / child & \textbf{person} / \textbf{ciald} \\
man / woman & \textbf{man} / \textbf{uuman} \\
mother-father & \textbf{muder} / \textbf{fader} \\
friend & \textbf{frend} \\
head / eye & \textbf{hed} / \textbf{ai} \\
mouth & \textbf{maut} \\
arm / leg & \textbf{arm} / \textbf{leg} \\
hand / foot & \textbf{hand} / \textbf{fut} \\
\end{longtable}
}

\subsubsection{4.4 Basic Adjectives \&
Colors}\label{basic-adjectives-colors}

{\def\LTcaptype{} % do not increment counter
\begin{longtable}[]{@{}ll@{}}
\toprule\noalign{}
English & MinEnglish \\
\midrule\noalign{}
\endhead
\bottomrule\noalign{}
\endlastfoot
big / small & \textbf{big} / \textbf{smol} \\
good / bad & \textbf{gud} / \textbf{bad} \\
new / old & \textbf{niu} / \textbf{ould} \\
red / blue & \textbf{red} / \textbf{bluu} \\
hot / cold & \textbf{hot} / \textbf{could} \\
happy / sad & \textbf{hapi} / \textbf{sad} \\
quick / slow & \textbf{cuic} / \textbf{slou} \\
green / black & \textbf{griin} / \textbf{blac} \\
\end{longtable}
}

\begin{center}\rule{0.5\linewidth}{0.5pt}\end{center}

\newpage

\subsection{5. Syntactic Reference
Guide}\label{syntactic-reference-guide}

\begin{Shaded}
\begin{Highlighting}[]
\NormalTok{NOUN:       \textless{}count\textgreater{}\textless{}noun\textgreater{}         3cat, *dog, \textasciitilde{}10person, 0eror}
\NormalTok{PROPER:     \textless{}count\textgreater{}\textless{}noun\textgreater{} \textless{}Name\textgreater{}  person Isabella, cuntri Spain}
\NormalTok{MATH:       \textless{}num\textgreater{}/\textless{}num\textgreater{}, \textless{}num\textgreater{}.\textless{}num\textgreater{}  1/2apel, 0.5bol}
\NormalTok{VERB:       (none)=present        sit, iit, run (default now)}
\NormalTok{            1 = present (attach)  h1laic, tinc}
\NormalTok{            * = any/all/habitual  *iit, *run}
\NormalTok{            Nd = future (+ optional)  1d gou, 2h start}
\NormalTok{            {-}Nd = past            {-}1d iit, {-}2h}
\NormalTok{            : = duration          :5Y studi (studying for 5 years)}
\NormalTok{ADJ:        \textless{}noun\textgreater{} \textless{}adj\textgreater{}          cat big, haus niu}
\NormalTok{INTENSITY:  \textgreater{}/\textless{} \textless{}adj\textgreater{}             cat \textgreater{}big (very big), cat \textless{}big (slightly)}
\NormalTok{MODAL:      \textless{}modal\textgreater{} \textless{}verb\textgreater{}        can run, shud iit, \textasciitilde{}can cum}
\NormalTok{NEG:        no \textless{}verb\textgreater{}             no laic, {-}1d no cum}
\NormalTok{QUEST:      ... ?                 u laic cat?}
\NormalTok{COMPARE:    mor / most \textless{}adj\textgreater{}      mor big, most big}
\NormalTok{APPROX:     \textasciitilde{}\textless{}value\textgreater{}              \textasciitilde{}5, \textasciitilde{}big}
\NormalTok{GENERIC:    *\textless{}noun\textgreater{}, *\textless{}verb\textgreater{}      *dog *laic *fuud}
\NormalTok{PASSIVE:    \textless{}Obj\textgreater{} bai \textless{}Subj\textgreater{}      bol kic bai i (ball kicked by me)}
\NormalTok{ATTACH:     \textless{}pronoun\textgreater{}\textless{}prefix\textgreater{}     h1laic, {-}1d, s1run}
\end{Highlighting}
\end{Shaded}

\subsubsection{Consonant Key}\label{consonant-key}

\begin{Shaded}
\begin{Highlighting}[]
\NormalTok{c = /k/ soft (cat, cum, bac)      k = /kʰ/ hard (king, kiip, kic)}
\NormalTok{ci / ce = /tʃ/ (ciald, cier, ceic) sh = /ʃ/   ng = /ŋ/   zh = /ʒ/}
\NormalTok{u = /w/ (uoter, uid)              i = /j/ (ies, iir)}
\NormalTok{t = /t/ (tinc, trii)              d = /d/ (de, dat)}
\end{Highlighting}
\end{Shaded}

\subsubsection{Pronouns}\label{pronouns}

\begin{Shaded}
\begin{Highlighting}[]
\NormalTok{i = I/me       u = you        h = he{-}him}
\NormalTok{s = she{-}her    t = it         dis = this      dat = that}

\NormalTok{Plural: 5i = we five, *u = you all, 3h = 3 males, *s = they{-}females, *t = they{-}things}
\NormalTok{Possess: i\textquotesingle{}s, u\textquotesingle{}s, h\textquotesingle{}s, s\textquotesingle{}, t\textquotesingle{}s}
\end{Highlighting}
\end{Shaded}

\begin{center}\rule{0.5\linewidth}{0.5pt}\end{center}

\newpage

\subsection{6. Comparative Corpus}\label{comparative-corpus}

\subsubsection{Simple \& Truths}\label{simple-truths}

\paragraph{Ex. 1}\label{ex.-1}

\begin{Shaded}
\begin{Highlighting}[]
\NormalTok{EN: The three cats are sitting on the big red mat. (46)}
\NormalTok{ME: 3cat sit on mat big red. (24)}
\NormalTok{BT: Three{-}cat sit on one{-}mat big red.}
\end{Highlighting}
\end{Shaded}

\begin{quote}
\textbf{Compression:} 47.8\% Δ

\textbf{Observation:} High compression achieved via removal of definite
articles, auxiliary verbs, and strict phonetic normalization.
\end{quote}

\paragraph{Ex. 2}\label{ex.-2}

\begin{Shaded}
\begin{Highlighting}[]
\NormalTok{EN: Dogs are loyal animals that love their owners. (46)}
\NormalTok{ME: *dog bi *animal loial dat *luv *h\textquotesingle{}s ouner. (42)}
\NormalTok{BT: Any{-}dog be any{-}animal loyal that any{-}love any{-}their owner.}
\end{Highlighting}
\end{Shaded}

\begin{quote}
\textbf{Compression:} 8.7\% Δ

\textbf{Observation:} Universal deference operators (\texttt{*})
substitute complex plural or aggregate noun constraints.
\end{quote}

\subsubsection{Relatives \& Times}\label{relatives-times}

\paragraph{Ex. 3}\label{ex.-3}

\begin{Shaded}
\begin{Highlighting}[]
\NormalTok{EN: Yesterday I bought two books and three apples at the store. (59)}
\NormalTok{ME: {-}1d bai 2buc an 3apel @ stor. (29)}
\NormalTok{BT: I one{-}day{-}ago buy two{-}book and three{-}apple at one{-}store.}
\end{Highlighting}
\end{Shaded}

\begin{center}\rule{0.5\linewidth}{0.5pt}\end{center}

\begin{quote}
\textbf{Compression:} 50.8\% Δ

\textbf{Observation:} Temporal operators (\texttt{-1d}, \texttt{1d},
\texttt{:5Y}) collapse complex tense architectures and prepositional
phrases.
\end{quote}

\newpage

\subsection{7. Comparative Information Density
Analysis}\label{comparative-information-density-analysis}

The transmission efficiency (information density) of MinEnglish was
modeled against a 35-sentence standard English corpus. Efficiency is
measured locally in byte-length (character count reduction).

\subsubsection{Top Compression Vectors
vs.~Deficits}\label{top-compression-vectors-vs.-deficits}

{\def\LTcaptype{} % do not increment counter
\begin{longtable}[]{@{}
  >{\raggedright\arraybackslash}p{(\linewidth - 4\tabcolsep) * \real{0.2500}}
  >{\raggedright\arraybackslash}p{(\linewidth - 4\tabcolsep) * \real{0.0714}}
  >{\raggedright\arraybackslash}p{(\linewidth - 4\tabcolsep) * \real{0.6786}}@{}}
\toprule\noalign{}
\begin{minipage}[b]{\linewidth}\raggedright
Highest Reductions
\end{minipage} & \begin{minipage}[b]{\linewidth}\raggedright
Char Δ
\end{minipage} & \begin{minipage}[b]{\linewidth}\raggedright
Variable Vector
\end{minipage} \\
\midrule\noalign{}
\endhead
\bottomrule\noalign{}
\endlastfoot
\textbf{Ex. 6:} Absolute Timestamp & \textbf{63.4\%} &
\texttt{2026-01-15\ 14:30} collapses 59 English alphabetical characters
into absolute numeric notation. \\
\textbf{Ex. 5:} Approximation & \textbf{58.6\%} &
\texttt{\textasciitilde{}50person} collapses multi-word approximation
phrases (``About fifty people''). \\
\textbf{Ex. 4, 12, 30:} Modality \& Tense &
\textbf{\textasciitilde45.0\%} & Prefixing operators (\texttt{1d},
\texttt{btuiin\ 12:00}) computationally squash auxiliary verb chains. \\
\end{longtable}
}

{\def\LTcaptype{} % do not increment counter
\begin{longtable}[]{@{}
  >{\raggedright\arraybackslash}p{(\linewidth - 4\tabcolsep) * \real{0.2183}}
  >{\raggedright\arraybackslash}p{(\linewidth - 4\tabcolsep) * \real{0.0634}}
  >{\raggedright\arraybackslash}p{(\linewidth - 4\tabcolsep) * \real{0.7183}}@{}}
\toprule\noalign{}
\begin{minipage}[b]{\linewidth}\raggedright
Lowest Reductions
\end{minipage} & \begin{minipage}[b]{\linewidth}\raggedright
Char Δ
\end{minipage} & \begin{minipage}[b]{\linewidth}\raggedright
Variable Vector
\end{minipage} \\
\midrule\noalign{}
\endhead
\bottomrule\noalign{}
\endlastfoot
\textbf{Ex. 18:} Lexical Description & \textbf{8.5\%} & Base nouns of
Latinate origin (\texttt{hospital}, \texttt{iuniversiti}) maintain
protracted phonetic transcriptions. \\
\textbf{Ex. 31:} Legal Domain & \textbf{20.0\%} & Domain-specific
transliterations (\texttt{diipendant}, \texttt{carj}) resist operator
optimization. \\
\end{longtable}
}

\textbf{Corpus Density Metric (n=35):}

\begin{itemize}
\tightlist
\item
  Total Standard English Bytes (Chars): \textbf{2633}
\item
  Total MinEnglish Bytes (Chars): \textbf{1711}
\item
  Average Information Density Gain (Length Reduction): \textbf{35.0\%}
\end{itemize}

\begin{center}\rule{0.5\linewidth}{0.5pt}\end{center}

\newpage

\subsection{8. Theoretical Limitations: The Human Biological
Factor}\label{theoretical-limitations-the-human-biological-factor}

The structural additions of previous linguistic protocols
(Literalization \texttt{\_}, Clause Anchors \texttt{co}/\texttt{oc}, and
Echo-Tagging) establish MinEnglish as a mathematically and
informationally optimized transmission system. It addresses standard
linguistic inefficiencies:

\begin{enumerate}
\def\labelenumi{\arabic{enumi}.}
\tightlist
\item
  \textbf{Short-Term Memory Bounds:} Clause anchors act as explicit
  stack managers for human working memory, permitting indefinite
  center-embedding.
\item
  \textbf{Transmission Integrity:} Echo-Tagging permits dynamic
  redundancy scaling to guarantee fidelity in high-loss environments
  (Shannon's Theorem).
\end{enumerate}

\subsubsection{8.1 Psychological Friction}\label{psychological-friction}

The primary identified limitation of MinEnglish exists externally to its
syntactic structure; the limitation is biological. Mammalian
communicative pathways evolved to facilitate social cohesion,
hierarchical navigation, and emotional subtext, often utilizing
ambiguity to reduce friction.

MinEnglish's computational design enforces an objective, algebraic
paradigm. The absolute requirement for temporal precision, numerical
exactness, and the prohibition of metaphorical obfuscation inherently
conflicts with natural human cognitive behavior. Widespread neurological
adoption of MinEnglish natively would likely prompt measurable
psychological resistance, as the protocol disables the linguistic
mechanisms humans rely on to process and mitigate sociological
complexity.

\subsubsection{8.2 Global Formalization
Standards}\label{global-formalization-standards}

As MinEnglish is designed to be a mathematically and informationally
optimized transmission system, it inherently relies on existing,
globally accepted objective standards to eliminate ambiguity. We
strongly recommend that all implementations, translations, and corpus
generations of MinEnglish strictly adhere to the following universal
formalisms:

\begin{enumerate}
\def\labelenumi{\arabic{enumi}.}
\tightlist
\item
  \textbf{Chronometry:} All temporal references should default to
  Universal Coordinated Time (UTC) or Universal Standard Time (UST)
  representations (e.g., \texttt{14:30Z}).
\item
  \textbf{Measurement:} Only the International System of Units (SI /
  Metric System) should be utilized for weights, measures, distances,
  and volumes.
\item
  \textbf{Mathematics:} Equations, variable definitions, and complex
  mathematical relationships should be expressed natively in LaTeX
  syntax to avoid ad-hoc lexical transcriptions.
\item
  \textbf{Dates:} Adherence strictly to the ISO 8601 standard
  (\texttt{YYYY-MM-DD}).
\end{enumerate}

\begin{center}\rule{0.5\linewidth}{0.5pt}\end{center}

\newpage

\subsection{9. References}\label{references}

\begin{itemize}
\tightlist
\item
  Shannon, C. E. (1948). A mathematical theory of communication.
  \emph{Bell System Technical Journal}, 27(3), 379--423.
\item
  Miller, G. A. (1956). The magical number seven, plus or minus two:
  Some limits on our capacity for processing information.
  \emph{Psychological Review}, 63(2), 81--97.
\item
  Chomsky, N. (1957). \emph{Syntactic Structures}. Mouton.
\item
  Zipf, G. K. (1949). \emph{Human Behavior and the Principle of Least
  Effort}. Addison-Wesley.
\item
  Ogden, C. K. (1930). \emph{Basic English: A General Introduction with
  Rules and Grammar}. Kegan Paul.
\item
  Zamenhof, L. L. (1887). \emph{Dr.~Esperanto's International Language}.
\item
  Pinker, S. (1994). \emph{The Language Instinct: How the Mind Creates
  Language}. William Morrow.
\item
  Sapir, E. (1921). \emph{Language: An Introduction to the Study of
  Speech}. Harcourt, Brace.
\end{itemize}

\begin{center}\rule{0.5\linewidth}{0.5pt}\end{center}

\newpage

\subsection{Appendix A: Full Comparative
Corpus}\label{appendix-a-full-comparative-corpus}

\subsubsection{Ex. 4}\label{ex.-4}

\begin{Shaded}
\begin{Highlighting}[]
\NormalTok{EN: Will you eat dinner with us tomorrow evening? (45)}
\NormalTok{ME: u1d iit diner uid *i? (22)}
\NormalTok{BT: You in{-}one{-}day eat one{-}dinner with all{-}of{-}us?}
\end{Highlighting}
\end{Shaded}

\begin{quote}
\textbf{Compression:} 51.1\% Δ

\textbf{Observation:} Temporal operators (\texttt{-1d}, \texttt{1d},
\texttt{:5Y}) collapse complex tense architectures and prepositional
phrases.
\end{quote}

\paragraph{Ex. 5}\label{ex.-5}

\begin{Shaded}
\begin{Highlighting}[]
\NormalTok{EN: About fifty people arrived at approximately three o\textquotesingle{}clock. (58)}
\NormalTok{ME: \textasciitilde{}50person ariv @ \textasciitilde{}15:00. (24)}
\NormalTok{BT: Approximately{-}fifty{-}person arrive at approximately{-}15:00.}
\end{Highlighting}
\end{Shaded}

\begin{quote}
\textbf{Compression:} 58.6\% Δ

\textbf{Observation:} Temporal operators (\texttt{-1d}, \texttt{1d},
\texttt{:5Y}) collapse complex tense architectures and prepositional
phrases.
\end{quote}

\paragraph{Ex. 6}\label{ex.-6}

\begin{Shaded}
\begin{Highlighting}[]
\NormalTok{EN: The meeting was held on January fifteenth, twenty twenty{-}six, at two thirty in the afternoon. (93)}
\NormalTok{ME: miiting 2026{-}01{-}15 14:30 bi held. (33)}
\NormalTok{BT: One{-}meeting 2026{-}01{-}15 14:30 be held.}
\end{Highlighting}
\end{Shaded}

\begin{quote}
\textbf{Compression:} 64.5\% Δ

\textbf{Observation:} Temporal operators (\texttt{-1d}, \texttt{1d},
\texttt{:5Y}) collapse complex tense architectures and prepositional
phrases.
\end{quote}

\paragraph{Ex. 7}\label{ex.-7}

\begin{Shaded}
\begin{Highlighting}[]
\NormalTok{EN: I have been studying English for five years but I still cannot speak fluently. (78)}
\NormalTok{ME: :5Y studi inglish but i !can spiic fluuent. (44)}
\NormalTok{BT: For{-}five{-}years study English but I negation{-}can speak fluent.}
\end{Highlighting}
\end{Shaded}

\begin{quote}
\textbf{Compression:} 42.3\% Δ

\textbf{Observation:} Temporal operators (\texttt{-1d}, \texttt{1d},
\texttt{:5Y}) collapse complex tense architectures and prepositional
phrases.
\end{quote}

\subsubsection{Complex Modifiers \&
Voice}\label{complex-modifiers-voice}

\paragraph{Ex. 8}\label{ex.-8}

\begin{Shaded}
\begin{Highlighting}[]
\NormalTok{EN: We are going to the extremely new restaurant tomorrow at seven in the evening. (78)}
\NormalTok{ME: *i1d gou tu fuud-haus >>niu @ 19:00. (39)}
\NormalTok{BT: All{-}of{-}us in{-}one{-}day go to one{-}food{-}house extremely{-}new at 19:00.}
\end{Highlighting}
\end{Shaded}

\begin{quote}
\textbf{Compression:} 51.3\% Δ

\textbf{Observation:} Temporal operators (\texttt{-1d}, \texttt{1d},
\texttt{:5Y}) collapse complex tense architectures and prepositional
phrases.
\end{quote}

\paragraph{Ex. 9}\label{ex.-9}

\begin{Shaded}
\begin{Highlighting}[]
\NormalTok{EN: Look out! The glass window was just broken by the angry man. (60)}
\NormalTok{ME: uoci!! uindou glas -0s breic bai man angri. (44)}
\NormalTok{BT: Watch! One{-}window glass just{-}now break by one{-}man angry.}
\end{Highlighting}
\end{Shaded}

\begin{quote}
\textbf{Compression:} 35.0\% Δ

\textbf{Observation:} Exclamatory prefixes (\texttt{!!}) enforce
immediate semantic parsing, bypassing syntactic buildup.
\end{quote}

\paragraph{Ex. 10}\label{ex.-10}

\begin{Shaded}
\begin{Highlighting}[]
\NormalTok{EN: She must give her book to him before she can start working. (59)}
\NormalTok{ME: s1must giv s\textquotesingle{}s buc tu h bifor s1can start tu uorc. (50)}
\NormalTok{BT: She{-}now{-}must give her book to he before she{-}now{-}can start [Object] one{-}work.}
\end{Highlighting}
\end{Shaded}

\begin{quote}
\textbf{Compression:} 15.3\% Δ

\textbf{Observation:} General lexical optimization and phonetic spelling
reduction.
\end{quote}

\subsubsection{Conditionals, Logic \&
Commands}\label{conditionals-logic-commands}

\paragraph{Ex. 11}\label{ex.-11}

\begin{Shaded}
\begin{Highlighting}[]
\NormalTok{EN: If you don\textquotesingle{}t stop running, you will fall and break your legs. (61)}
\NormalTok{ME: if u1!stop tu run, u1fol an breic tu u\textquotesingle{}s 2leg. (48)}
\NormalTok{BT: If you{-}now{-}negation{-}stop [Obj] run, you{-}now{-}fall and break [Obj] your two{-}leg.}
\end{Highlighting}
\end{Shaded}

\begin{quote}
\textbf{Compression:} 21.3\% Δ

\textbf{Observation:} General lexical optimization and phonetic spelling
reduction.
\end{quote}

\paragraph{Ex. 12}\label{ex.-12}

\begin{Shaded}
\begin{Highlighting}[]
\NormalTok{EN: Every morning I drink two cups of coffee and eat one piece of toast before work. (80)}
\NormalTok{ME: *morning i*drinc 2cup cofi an *iit piis toust bifor uorc. (58)}
\NormalTok{BT: Any{-}morning I{-}habitually{-}drink two{-}cup coffee and habitually{-}eat one{-}piece toast before work.}
\end{Highlighting}
\end{Shaded}

\begin{quote}
\textbf{Compression:} 27.5\% Δ

\textbf{Observation:} Universal deference operators (\texttt{*})
substitute complex plural or aggregate noun constraints.
\end{quote}

\paragraph{Ex. 13}\label{ex.-13}

\begin{Shaded}
\begin{Highlighting}[]
\NormalTok{EN: Take three eggs, add two cups of flour, and mix everything together for ten minutes. (84)}
\NormalTok{ME: teic 3eg, ad 2cup flaur, an mics *ting for 10m. (48)}
\NormalTok{BT: Take three{-}egg, add two{-}cup flour, and mix all{-}thing for ten{-}minutes.}
\end{Highlighting}
\Shaded}

\begin{quote}
\textbf{Compression:} 42.9\% Δ

\textbf{Observation:} Universal deference operators (\texttt{*})
substitute complex plural or aggregate noun constraints.
\end{quote}

\paragraph{Ex. 14}\label{ex.-14}

\begin{Shaded}
\begin{Highlighting}[]
\NormalTok{EN: Please give me two glasses of water, I am very thirsty. (55)}
\NormalTok{ME: \textasciitilde{}u giv i 2glas uoter, i bi \textgreater{}tersti. (35)}
\NormalTok{BT: Polite{-}you{-}give I two{-}glass water, I be very{-}thirsty.}
\end{Highlighting}
\Shaded}

\begin{quote}
\textbf{Compression:} 36.4\% Δ

\textbf{Observation:} Approximation operators
(\texttt{\textasciitilde{}}) drastically reduce multi-word phrasings
like ``about'' or ``approximately''.
\end{quote}

\subsubsection{Dialogue \& Description}\label{dialogue-description}

\paragraph{Ex. 15}\label{ex.-15}

\begin{Shaded}
\begin{Highlighting}[]
\NormalTok{EN: "How many children do you have?" "I have two boys and one girl." (64)}
\NormalTok{ME: "haum *ciald u hav?" "i hav 2boi an gerl." (42)}
\NormalTok{BT: "How{-}many any{-}child you have?" "I have two{-}boy and one{-}girl."}
\end{Highlighting}
\Shaded}

\begin{quote}
\textbf{Compression:} 35.9\% Δ

\textbf{Observation:} Interrogative prefixes condense investigative
clauses into single tokens.
\end{quote}

\paragraph{Ex. 16}\label{ex.-16}

\begin{Shaded}
\begin{Highlighting}[]
\NormalTok{EN: Last year, she traveled to five different countries and learned three new languages because she wanted to understand the world better. (134)}
\NormalTok{ME: s{-}1Y travel tu 5cuntri diferent an lern 3languij niu cos s{-}1Y uant understand uerld mor gud. (92)}
\NormalTok{BT: She one{-}year{-}ago travel to five{-}country different and learn three{-}language new because she one{-}year{-}ago want understand one{-}world more good.}
\end{Highlighting}
\Shaded}

\begin{quote}
\textbf{Compression:} 31.3\% Δ

\textbf{Observation:} Temporal operators (\texttt{-1d}, \texttt{1d},
\texttt{:5Y}) collapse complex tense architectures and prepositional
phrases.
\end{quote}

\paragraph{Ex. 17}\label{ex.-17}

\begin{Shaded}
\begin{Highlighting}[]
\NormalTok{EN: I might go to the party, but she loves him and he does not love her. (68)}
\NormalTok{ME: i1d \textasciitilde{}can gou tu parti, but s1luv h an h !luv s. (50)}
\NormalTok{BT: I in{-}future might go to one{-}party, but she{-}now{-}love he and he negation{-}love she.}
\end{Highlighting}
\Shaded}

\begin{quote}
\textbf{Compression:} 26.5\% Δ

\textbf{Observation:} Temporal operators (\texttt{-1d}, \texttt{1d},
\texttt{:5Y}) collapse complex tense architectures and prepositional
phrases.
\end{quote}

\paragraph{Ex. 18}\label{ex.-18}

\begin{Shaded}
\begin{Highlighting}[]
\NormalTok{EN: My brother works at a hospital and my sister studies at the university. (71)}
\NormalTok{ME: i\textquotesingle{}s bruder *uorc @ hospital an i\textquotesingle{}s sister *studi @ iuniversiti. (63)}
\NormalTok{BT: My brother any{-}work at one{-}hospital and my sister any{-}study at one{-}university.}
\end{Highlighting}
\Shaded}

\begin{quote}
\emph{Low compression --- long Latin words maintain their length.}
11.3\% Δ

\textbf{Observation:} Universal deference operators (\texttt{*})
substitute complex plural or aggregate noun constraints.
\end{quote}

\paragraph{Ex. 19}\label{ex.-19}

\begin{Shaded}
\begin{Highlighting}[]
\NormalTok{EN: The children were playing in the park when it started raining heavily. (70)}
\NormalTok{ME: *ciald in parc uen t{-}0s start rein \textgreater{}hevi. (42)}
\NormalTok{BT: Any{-}child just{-}now in(verb) one{-}park when it{-}just{-}now start rain very{-}heavily.}
\end{Highlighting}
\Shaded}

\begin{quote}
\textbf{Compression:} 40.0\% Δ

\textbf{Observation:} Intensity operators (\texttt{\textgreater{}},
\texttt{\textgreater{}\textgreater{}}) eliminate lengthy adverbial
modifiers.
\end{quote}

\paragraph{Ex. 20}\label{ex.-20}

\begin{Shaded}
\begin{Highlighting}[]
\NormalTok{EN: How much does this car cost? It is too expensive for me. (56)}
\NormalTok{ME: haum dis car cost? t bi \textgreater{}\textgreater{}ecspensiv for i. (43)}
\NormalTok{BT: How{-}much this car cost? It be extremely{-}expensive for I.}
\end{Highlighting}
\Shaded}

\begin{quote}
\textbf{Compression:} 23.2\% Δ

\textbf{Observation:} Intensity operators (\texttt{\textgreater{}},
\texttt{\textgreater{}\textgreater{}}) eliminate lengthy adverbial
modifiers.
\end{quote}

\subsubsection{Advanced Domains (Sci-Fi, Philosophy,
Business)}\label{advanced-domains-sci-fi-philosophy-business}

\paragraph{Ex. 21}\label{ex.-21}

\begin{Shaded}
\begin{Highlighting}[]
\NormalTok{EN: The spaceship launched into orbit after the countdown reached zero. (67)}
\NormalTok{ME: speisship -0s lonci in orbit after caunt{-}daun hit 0. (52)}
\NormalTok{BT: One{-}spaceship just{-}now launch in one{-}orbit after one{-}count{-}down hit zero.}
\end{Highlighting}
\Shaded}

\begin{quote}
\textbf{Compression:} 28.4\% Δ

\textbf{Observation:} General lexical optimization and phonetic spelling
reduction.
\end{quote}

\paragraph{Ex. 22}\label{ex.-22}

\begin{Shaded}
\begin{Highlighting}[]
\NormalTok{EN: To be or not to be, that is the question we must all answer eventually. (71)}
\NormalTok{ME: bi or !bi, dat bi cuesciun *i must anser in taim. (51)}
\NormalTok{BT: Be or not{-}be, that be one{-}question all{-}of{-}us must answer in future{-}time.}
\end{Highlighting}
\end{Shaded}

\begin{quote}
\textbf{Compression:} 26.8\% Δ

\textbf{Observation:} Universal deference operators (\texttt{*})
substitute complex plural or aggregate noun constraints.
\end{quote}

\paragraph{Ex. 23}\label{ex.-23}

\begin{Shaded}
\begin{Highlighting}[]
\NormalTok{EN: Our company needs to optimize its supply chain to maximize profit margins next year. (84)}
\NormalTok{ME: *i\textquotesingle{}s compani niid optimaiz t\textquotesingle{}s suplai{-}cein tu macsimaiz *marjin{-}profit 1Y. (74)}
\NormalTok{BT: All{-}of{-}us{-}possessive company need optimize its supply{-}chain to maximize any{-}margin{-}profit next{-}year.}
\end{Highlighting}
\end{Shaded}

\begin{quote}
\textbf{Compression:} 11.9\% Δ

\textbf{Observation:} Temporal operators (\texttt{-1d}, \texttt{1d},
\texttt{:5Y}) collapse complex tense architectures and prepositional
phrases.
\end{quote}

\paragraph{Ex. 24}\label{ex.-24}

\begin{Shaded}
\begin{Highlighting}[]
\NormalTok{EN: The president announced a new tax policy that will affect middle class families. (80)}
\NormalTok{ME: prezident anauns tacs{-}polisi niu dat 1d afect *famili clas{-}midl. (65)}
\NormalTok{BT: One{-}president just{-}now announce one{-}tax{-}policy new that in{-}future affect any{-}family class{-}middle.}
\end{Highlighting}
\end{Shaded}

\begin{quote}
\textbf{Compression:} 18.8\% Δ

\textbf{Observation:} Temporal operators (\texttt{-1d}, \texttt{1d},
\texttt{:5Y}) collapse complex tense architectures and prepositional
phrases.
\end{quote}

\paragraph{Ex. 25}\label{ex.-25}

\begin{Shaded}
\begin{Highlighting}[]
\NormalTok{EN: The doctor prescribed antibiotics to treat the patient\textquotesingle{}s severe bacterial infection. (84)}
\NormalTok{ME: doctor prescraib tu antibiotik tu triit bacteriel infecshun \textgreater{}siirius peishent\textquotesingle{}s. (80)}
\NormalTok{BT: One{-}doctor just{-}now prescribe [Obj] one{-}antibiotic to treat one{-}bacterial{-}infection very{-}serious one{-}patient\textquotesingle{}s.}
\end{Highlighting}
\end{Shaded}

\begin{quote}
\textbf{Compression:} 4.8\% Δ

\textbf{Observation:} Intensity operators (\texttt{\textgreater{}},
\texttt{\textgreater{}\textgreater{}}) eliminate lengthy adverbial
modifiers.
\end{quote}

\paragraph{Ex. 26}\label{ex.-26}

\begin{Shaded}
\begin{Highlighting}[]
\NormalTok{EN: The three bedroom house with a large backyard is currently on the market for rent. (82)}
\NormalTok{ME: haus 3bedruum uid baciard big in marcet for rent nau. (53)}
\NormalTok{BT: One{-}house three{-}bedroom with one{-}backyard big now{-}in one{-}market for rent now.}
\end{Highlighting}
\end{Shaded}

\begin{quote}
\textbf{Compression:} 35.4\% Δ

\textbf{Observation:} General lexical optimization and phonetic spelling
reduction.
\end{quote}

\paragraph{Ex. 27}\label{ex.-27}

\begin{Shaded}
\begin{Highlighting}[]
\NormalTok{EN: If the server crashes, the backup system will automatically restore the database. (81)}
\NormalTok{ME: if server craish, backup sistem 1d otoristor daatabeis. (55)}
\NormalTok{BT: If one{-}server crash, one{-}backup{-}system in{-}future auto{-}restore one{-}database.}
\end{Highlighting}
\end{Shaded}

\begin{quote}
\textbf{Compression:} 32.1\% Δ

\textbf{Observation:} Temporal operators (\texttt{-1d}, \texttt{1d},
\texttt{:5Y}) collapse complex tense architectures and prepositional
phrases.
\end{quote}

\paragraph{Ex. 28}\label{ex.-28}

\begin{Shaded}
\begin{Highlighting}[]
\NormalTok{EN: Global warming is melting the polar ice caps faster than scientists predicted. (78)}
\NormalTok{ME: gloubal uorming melt *cap{-}ais poular \textgreater{}fast dan *saientist {-}1d pridiict. (71)}
\NormalTok{BT: Global{-}warming now{-}melt any{-}cap{-}ice polar more{-}fast than any{-}scientist past{-}predict.}
\end{Highlighting}
\end{Shaded}

\begin{quote}
\textbf{Compression:} 9.0\% Δ

\textbf{Observation:} Temporal operators (\texttt{-1d}, \texttt{1d},
\texttt{:5Y}) collapse complex tense architectures and prepositional
phrases.
\end{quote}

\paragraph{Ex. 29}\label{ex.-29}

\begin{Shaded}
\begin{Highlighting}[]
\NormalTok{EN: The wind whispered through the ancient trees as the silver moon rose in the sky. (80)}
\NormalTok{ME: uind uisper tru *trii \textgreater{}ould az muun silver raiz in scai. (57)}
\NormalTok{BT: One{-}wind just{-}now whisper through any{-}tree very{-}old as one{-}moon silver just{-}now rise in one{-}sky.}
\end{Highlighting}
\end{Shaded}

\begin{quote}
\textbf{Compression:} 28.7\% Δ

\textbf{Observation:} Intensity operators (\texttt{\textgreater{}},
\texttt{\textgreater{}\textgreater{}}) eliminate lengthy adverbial
modifiers.
\end{quote}

\paragraph{Logistics}\label{logistics}

\paragraph{Ex. 30}\label{ex.-30}

\begin{Shaded}
\begin{Highlighting}[]
\NormalTok{EN: The delivery truck will arrive with your package between noon and two o\textquotesingle{}clock. (78)}
\NormalTok{ME: diliveri truc 1d ariv uid u\textquotesingle{}s pacij btuiin 12:00 an 14:00. (58)}
\NormalTok{BT: One{-}delivery{-}truck in{-}future arrive with your package between 12:00 and 14:00.}
\end{Highlighting}
\end{Shaded}

\begin{quote}
\textbf{Compression:} 25.6\% Δ

\textbf{Observation:} Temporal operators (\texttt{-1d}, \texttt{1d},
\texttt{:5Y}) collapse complex tense architectures and prepositional
phrases.
\end{quote}

\paragraph{Ex. 31}\label{ex.-31}

\begin{Shaded}
\begin{Highlighting}[]
\NormalTok{EN: The defendant formally pleaded not guilty to all charges presented by the court. (80)}
\NormalTok{ME: diifendant formal pliid !gilti tu *ciarj prizent bai cort. (59)}
\NormalTok{BT: One{-}defendant just{-}now formal plead not guilty to any{-}charge present by one{-}court.}
\end{Highlighting}
\end{Shaded}

\begin{quote}
\textbf{Compression:} 21.3\% Δ

\textbf{Observation:} Universal deference operators (\texttt{*})
substitute complex plural or aggregate noun constraints.
\end{quote}

\paragraph{Ex. 32}\label{ex.-32}

\begin{Shaded}
\begin{Highlighting}[]
\NormalTok{EN: Boil the water, add a pinch of salt, and stir constantly until the sauce thickens. (82)}
\NormalTok{ME: boil tu uoter, ad tu pince solt, an *ster until sos \textgreater{}tic. (58)}
\NormalTok{BT: Boil [Obj] one{-}water, add [Obj] one{-}pinch salt, and any{-}stir until one{-}sauce very{-}thick.}
\end{Highlighting}
\end{Shaded}

\begin{quote}
\textbf{Compression:} 29.3\% Δ

\textbf{Observation:} Intensity operators (\texttt{\textgreater{}},
\texttt{\textgreater{}\textgreater{}}) eliminate lengthy adverbial
modifiers.
\end{quote}

\paragraph{Ex. 33}\label{ex.-33}

\begin{Shaded}
\begin{Highlighting}[]
\NormalTok{EN: I have never felt this way about anyone before, you changed my entire life. (75)}
\NormalTok{ME: {-}1d *!fiil dis ue abaut *person bifor, u {-}1d ceinj i\textquotesingle{}s laif \textgreater{}hol. (67)}
\NormalTok{BT: I{-}just{-}now{-}yesterday any{-}negation{-}feel this way about any{-}person before, you yesterday change [Obj] my life very{-}whole.}
\end{Highlighting}
\end{Shaded}

\begin{quote}
\textbf{Compression:} 10.7\% Δ

\textbf{Observation:} Temporal operators (\texttt{-1d}, \texttt{1d},
\texttt{:5Y}) collapse complex tense architectures and prepositional
phrases.
\end{quote}

\paragraph{Ex. 34}\label{ex.-34}

\begin{Shaded}
\begin{Highlighting}[]
\NormalTok{EN: Go straight for two miles, then turn left at the second traffic light you see. (78)}
\NormalTok{ME: gou streit for 2mail, den tern left @ trapic lait tuu u1sii. (60)}
\NormalTok{BT: Go straight for two{-}mile, then turn left at one{-}traffic{-}light two you{-}now{-}see.}
\end{Highlighting}
\end{Shaded}

\begin{quote}
\textbf{Compression:} 23.1\% Δ

\textbf{Observation:} General lexical optimization and phonetic spelling
reduction.
\end{quote}

\paragraph{Ex. 35}\label{ex.-35}

\begin{Shaded}
\begin{Highlighting}[]
\NormalTok{EN: Hyperinflation destroys purchasing power and severely halts economic development. (81)}
\NormalTok{ME: \textgreater{}\textgreater{}haiperinflashun distroi tu pauer{-}bai an \textgreater{}siirius holt tu divelopment economik. (80)}
\NormalTok{BT: Extremely{-}hyperinflation now{-}destroy [Obj] power{-}buy and very{-}serious halt [Obj] development economic.}
\end{Highlighting}
\end{Shaded}

\begin{quote}
\textbf{Compression:} 1.2\% Δ

\textbf{Observation:} Low compression due to retention of exact
international academic vocabulary.
\end{quote}

\newpage

\subsection{Appendix B: License}\label{appendix-b-license}

\begin{Shaded}
\begin{Highlighting}[]
\NormalTok{MIT License}

\NormalTok{Copyright (c) 2026 Dan Micsa, PhD}

\NormalTok{Permission is hereby granted, free of charge, to any person obtaining a copy}
\NormalTok{of this language specification and associated documentation files (the}
\NormalTok{"Specification"), to deal in the Specification without restriction, including}
\NormalTok{without limitation the rights to use, copy, modify, merge, publish,}
\NormalTok{distribute, sublicense, and{-}or sell copies of the Specification, and to}
\NormalTok{permit persons to whom the Specification is furnished to do so, subject to}
\NormalTok{the following conditions:}

\NormalTok{The above copyright notice and this permission notice shall be included in}
\NormalTok{all copies or substantial portions of the Specification.}

\NormalTok{THE SPECIFICATION IS PROVIDED "AS IS", WITHOUT WARRANTY OF ANY KIND, EXPRESS}
\NormalTok{OR IMPLIED, INCLUDING BUT NOT LIMITED TO THE WARRANTIES OF MERCHANTABILITY,}
\NormalTok{FITNESS FOR A PARTICULAR PURPOSE AND NONINFRINGEMENT. IN NO EVENT SHALL THE}
\NormalTok{AUTHORS OR COPYRIGHT HOLDERS BE LIABLE FOR ANY CLAIM, DAMAGES OR OTHER}
\NormalTok{LIABILITY, WHETHER IN AN ACTION OF CONTRACT, TORT OR OTHERWISE, ARISING}
\NormalTok{FROM, OUT OF OR IN CONNECTION WITH THE SPECIFICATION OR THE USE OR OTHER}
\NormalTok{DEALINGS IN THE SPECIFICATION.}
\end{Highlighting}
\end{Shaded}
